\subsection{Spectral Radius and Matrix Powers}

Throughout this subsection, $\mathbb{F}$ is a field.

\begin{theorem} \label{thm:ranks-of-matrix-powers-stabilize} \hlt{Ranks of Matrix Powers Stabilize}\\
Let $n\in\mathbb{Z}_{>0}$ and $A\in\operatorname{Mat}(n,n,\mathbb{F})$. Then:
\begin{enumerate}[label=(\roman*)]
\setlength{\itemsep}{0pt}
\item $(\forall\,k\in\mathbb{Z}_{>0})\,\operatorname{rank}(A^k)\geq\operatorname{rank}(A^{k+1})$;
\item if $k\in\mathbb{Z}_{>0}$ and $\operatorname{rank}(A^k)=\operatorname{rank}(A^{k+1})$, then $(\forall\,p\in\mathbb{Z}_{>0})\,\bigl[\operatorname{col}(A^k)=\operatorname{col}(A^{k+p})\ \land\ \operatorname{rank}(A^k)=\operatorname{rank}(A^{k+p})\bigr]$;
\item if $A$ is not invertible, then $\min\{q\in\mathbb{Z}_{>0}:\operatorname{rank}(A^q)=\operatorname{rank}(A^{q+1})\}\in J_n$.
\end{enumerate}
\begin{proof}
\begin{enumerate}[label=(\roman*)]
\setlength{\itemsep}{0pt}
\item Fix $k\in\mathbb{Z}_{>0}$. Then $\operatorname{col}(A^{k+1})\subseteq\operatorname{col}(A^k)$ by $A^{k+1}=A^kA$, and the rank inequality holds by taking dimensions.
\item Fix $k\in\mathbb{Z}_{>0}$ with $\operatorname{rank}(A^k)=\operatorname{rank}(A^{k+1})$, and set $U\defeq\operatorname{col}(A^k)$.\\
The identities $U=\operatorname{col}(A^{k+1})$ and $A(U)=U$ hold by part (i), equality of dimensions, and $AA^k=A^{k+1}$.\\
$(\forall\,p\in\mathbb{Z}_{>0})\,\operatorname{col}(A^{k+p})=U$ by induction; the rank identities hold by taking dimensions.
\item $\operatorname{rank}(A)\leq n-1$ by Corollary \ref{cor:square-matrix-invertible-iff-full-rank} and noninvertibility.\\
If $(\forall\,q\in J_n)\,\operatorname{rank}(A^q)>\operatorname{rank}(A^{q+1})$, then $\operatorname{rank}(A^{n+1})\leq\operatorname{rank}(A)-n\leq-1$, contradicting nonnegativity of rank.\\
The relation $(\exists\,q\in J_n)\,\operatorname{rank}(A^q)=\operatorname{rank}(A^{q+1})$ holds by this contradiction.\\
The minimum in (iii) belongs to $J_n$ by this witness and well-ordering.
\end{enumerate}
\hlt{Structural Commentary}: Nested column spaces and dimension equality create an invariant image; finite dimension forces stabilization.\\
\hlt{Key Insight}: Each strict rank drop consumes at least one dimension.
\end{proof}
\end{theorem}

\begin{definition} \label{def:matrix-index} \hlt{Index of a Matrix}\\
Let $n\in\mathbb{Z}_{>0}$ and $A\in\operatorname{Mat}(n,n,\mathbb{F})$.
\begin{enumerate}[label=(\roman*)]
\setlength{\itemsep}{0pt}
\item If $A$ is invertible, its \hlt{index} is $0$.
\item If $A$ is not invertible, its \hlt{index} is $\min\{q\in\mathbb{Z}_{>0}:\operatorname{rank}(A^q)=\operatorname{rank}(A^{q+1})\}$.
\end{enumerate}
\end{definition}

\begin{definition} \label{def:spectral-radius-and-convergent-matrix} \hlt{Spectral Radius and Convergent Matrix}\\
Let $n\in\mathbb{Z}_{>0}$ and $A\in\operatorname{Mat}(n,n,\mathbb{C})$.
\begin{enumerate}[label=(\roman*)]
\setlength{\itemsep}{0pt}
\item The \hlt{spectral radius} is $\rho(A)\defeq\max\limits_{\lambda\in\operatorname{spec}(A)}\abs{\lambda}$.
\item The notation $A^p\to0$ means $(\forall\,(i,j)\in J_n^2)\,\lim\limits_{p\to\infty}(A^p)_{ij}=0$.
\end{enumerate}
\end{definition}

\begin{proposition} \label{prop:nilpotent-matrix-basic-invariants} \hlt{Basic Invariants of Nilpotents}\\
Let $n\in\mathbb{Z}_{>0}$ and $N\in\operatorname{Mat}(n,n,\mathbb{F})$. The following are equivalent:
\begin{enumerate}[label=(\roman*)]
\setlength{\itemsep}{0pt}
\item $N$ is nilpotent;
\item the characteristic polynomial of $N$ is $z^n$;
\item $N^n=0$.
\end{enumerate}
Under these equivalent conditions, $\operatorname{tr}(N)=\det(N)=0$ and, if $\mathbb{F}=\mathbb{C}$, $\rho(N)=0$.
\begin{proof}
\begin{enumerate}[label=(\roman*)]
\setlength{\itemsep}{0pt}
\item[$(i)$]$\Rightarrow (ii)$ By nilpotency, choose $r\in\mathbb{Z}_{>0}$ satisfying $N^r=0$.\\
Every eigenvalue $\lambda$ over an algebraic closure equals $0$ by $\lambda^r\mathbf{v}=N^r\mathbf{v}=\mathbf{0}$ for an associated $\mathbf{v}\neq\mathbf{0}$.\\
$p_N(z)=z^n$ by complete splitting and Theorem \ref{thm:matrix-characteristic-polynomial-basic-coefficients}(i). \item[$(ii)$]$\Rightarrow (iii)$ $p_N(N)=N^n=0$ by Theorem \ref{thm:cayley-hamilton-matrix}. \item[$(iii)$]$\Rightarrow (i)$ Matrix $N$ is nilpotent by $N^n=0$.
\end{enumerate}
Assume the equivalent conditions. The identities $\operatorname{tr}(N)=\det(N)=0$ hold by $p_N(z)=z^n$ and Theorem \ref{thm:matrix-characteristic-polynomial-basic-coefficients}(i).\\
If $\mathbb{F}=\mathbb{C}$, then $\operatorname{spec}(N)=\{0\}$ by Theorem \ref{thm:matrix-characteristic-polynomial-basic-coefficients}(ii), and $\rho(N)=0$ by the spectral-radius formula.\\
\hlt{Structural Commentary}: Eigenvalue collapse identifies $p_N$; Cayley--Hamilton and its coefficients recover the remaining invariants.\\
\hlt{Key Insight}: Vanishing eigenvalues precede the uniform exponent $N^n=0$.
\end{proof}
\end{proposition}

\begin{theorem} \label{thm:matrix-powers-converge-to-zero-spectral-radius} \hlt{Matrix Powers Converge to Zero}\\
Let $n\in\mathbb{Z}_{>0}$ and $A\in\operatorname{Mat}(n,n,\mathbb{C})$. Then $A^p\to0$ if and only if $\rho(A)<1$.
\begin{proof}
By Theorem \ref{thm:complex-jordan-form}, choose $S\in\operatorname{Mat}(n,n,\mathbb{C})^\times$ and a Jordan matrix $J$ satisfying $A=SJS^{-1}$.\\
$A^p\to0\Longleftrightarrow J^p\to0$ by $A^p=SJ^pS^{-1}$, $J^p=S^{-1}A^pS$, and finite sums.\\
The convergence $J^p\to0$ is equivalent to convergence of each Jordan block power by the block-diagonal structure of $J$.
\begin{enumerate}[label=\roman*.]
\setlength{\itemsep}{0pt}
\item[$(\Rightarrow)$] $\rho(A)=\max\limits_{\lambda\in\operatorname{spec}(A)}\abs{\lambda}<1$ by the block reduction and Proposition \ref{prop:powers-of-jordan-block}(ii).
\item[$(\Leftarrow)$] Each Jordan block power tends to $0$ by $\rho(A)<1$ and Proposition \ref{prop:powers-of-jordan-block}(ii). $A^p\to0$ by the two reductions above.
\end{enumerate}
\hlt{Structural Commentary}: Similarity transfers entrywise convergence to block powers, which reduce it to eigenvalue moduli.\\
\hlt{Key Insight}: Geometric decay dominates each fixed polynomial superdiagonal factor.
\end{proof}
\end{theorem}

Figure \ref{fig:jordan-spectral-radius-convergence} depicts the Jordan-block threshold at the unit circle.

\begin{illustrationfigure}[H]
  \centering
  \begin{tikzpicture}[
  x=0.86cm,
  y=0.86cm,
  font=\small,
  box/.style={draw=blue!65!black, fill=blue!7, rounded corners=2pt, minimum width=2.45cm, minimum height=0.68cm, align=center},
  good/.style={draw=green!45!black, fill=green!8, rounded corners=2pt, minimum width=2.0cm, minimum height=0.62cm, align=center},
  bad/.style={draw=red!60!black, fill=red!7, rounded corners=2pt, minimum width=2.0cm, minimum height=0.62cm, align=center},
  arrow/.style={-{Stealth[length=2mm]}, thick, gray!75}
]
  \node[box] (jp) at (0,0) {$J_k(\lambda)^p$};
  \node[box] (formula) at (3.7,0) {$\binom{p}{r}\lambda^{p-r}$ terms};
  \node[good] (inside) at (7.2,0.75) {$\abs{\lambda}<1$\\decays};
  \node[bad] (outside) at (7.2,-0.75) {$\abs{\lambda}\geq1$\\does not decay};

  \draw[arrow] (jp) -- (formula);
  \draw[arrow] (formula.east) -- (inside.west);
  \draw[arrow] (formula.east) -- (outside.west);

  \draw[gray!60] (10,0) circle (0.9);
  \fill[gray!65] (10,0) circle (1.2pt) node[below, black] {\scriptsize $0$};
  \fill[green!60!black] (9.65,0.35) circle (2pt) node[above left, black] {\scriptsize inside};
  \fill[red!65!black] (10.9,0) circle (2pt) node[right, black] {\scriptsize on $\abs{\lambda}=1$};
\end{tikzpicture}

  \caption{Jordan block powers decay for $\abs{\lambda}<1$, but not on or outside the gray unit circle.}
  \label{fig:jordan-spectral-radius-convergence}
\end{illustrationfigure}

\begin{definition} \label{def:power-bounded-matrix} \hlt{Power Bounded Matrix}\\
Let $n\in\mathbb{Z}_{>0}$ and $A\in\operatorname{Mat}(n,n,\mathbb{C})$. The matrix $A$ is \hlt{power bounded} if $(\forall\,(i,j)\in J_n^2)\,\sup\limits_{p\in\mathbb{Z}_{>0}}\abs{(A^p)_{ij}}<\infty$.
\end{definition}

\begin{theorem} \label{thm:matrix-power-bounded-jordan-criterion} \hlt{Power Bounded Jordan Criterion}\\
Let $n\in\mathbb{Z}_{>0}$ and $A\in\operatorname{Mat}(n,n,\mathbb{C})$. The following are equivalent:
\begin{enumerate}[label=(\roman*)]
\setlength{\itemsep}{0pt}
\item $A$ is power bounded.
\item $\rho(A)\leq1$, and Jordan blocks associated with eigenvalues of modulus $1$ have size $1$.
\item $\rho(A)\leq1$, and $(\forall\,\lambda\in\operatorname{spec}(A))\,\bigl(\abs{\lambda}=1\Longrightarrow\dim E_{\lambda}(A)=m_A(\lambda)\bigr)$.
\end{enumerate}
\begin{proof}
By Theorem \ref{thm:complex-jordan-form}, choose $S\in\operatorname{Mat}(n,n,\mathbb{C})^\times$ and a Jordan matrix $J$ satisfying $A=SJS^{-1}$.\\
Power boundedness of $A$ and $J$ is equivalent by $A^p=SJ^pS^{-1}$, $J^p=S^{-1}A^pS$, and entrywise fixed finite linear combinations.
\begin{enumerate}[label=(\roman*)]
\setlength{\itemsep}{0pt}
\item[$(i)$]$\Rightarrow (ii)$ The matrix $J$ is power bounded by (i) and the fixed-similarity equivalence.\\
A block with $\abs{\lambda}>1$ is excluded by its unbounded diagonal sequence $(\lambda^p)_{p\geq1}$.\\
If $\abs{\lambda}=1$ and $k\geq2$, then $J_k(\lambda)^p$ is unbounded by Proposition \ref{prop:powers-of-jordan-block}(i) and $\abs{p\lambda^{p-1}}=p$. Condition (ii) holds by these exclusions. \item[$(ii)$]$\Rightarrow (iii)$ Fix $\lambda\in\operatorname{spec}(A)$ with $\abs{\lambda}=1$.\\
The dimension $\dim E_{\lambda}(A)$ is the number of Jordan blocks at $\lambda$ by Corollary \ref{cor:jordan-blocks-diagonalizability}(i).\\
The multiplicity $m_A(\lambda)$ is the sum of their sizes by Theorem \ref{thm:similarity-preserves-characteristic-polynomial}(ii) and Proposition \ref{prop:upper-triangular-characteristic-polynomial}(iii).\\
The desired equality holds by these two counting identities and the size-one condition in (ii). \item[$(iii)$]$\Rightarrow (i)$ All unit-circle blocks have size $1$ by (iii), the preceding counts, and positive integral block sizes.\\
Blocks with $\abs{\lambda}<1$ converge to $0$ by Proposition \ref{prop:powers-of-jordan-block}(ii), while the size-one unit-circle blocks are bounded by $\abs{\lambda^p}=1$.\\
The matrix $J$ is power bounded by $\rho(A)\leq1$, these block bounds, and its block-diagonal structure.\\
The matrix $A$ is power bounded by the fixed-similarity equivalence.
\end{enumerate}
\hlt{Structural Commentary}: Jordan form separates exponential and boundary polynomial growth; block counts detect nilpotent parts.\\
\hlt{Key Insight}: A unit-circle block is bounded exactly when its size is one.
\end{proof}
\end{theorem}
